\section{Polynomial commitments from SPIN}
\label{sec:spin-pcs}

We use SPIN as the row encoder in a Brakedown-style polynomial commitment
scheme~\cite{brakedown23}. This application tests whether fast ordinary
encoding translates into fast commitment and opening. The implementation
has two principal changes: it instantiates the code with SPIN, and it
organizes the computation around packed binary data. At $512$\,MiB of
input, the resulting prover commits and opens in $520$\,ms on one core.

\paragraph{Construction.}
Let a Boolean evaluation table define a multilinear polynomial over
$\mathbb F:=\F_{2^{128}}$. A public index mapping arranges its $RK$ bits
into $R$ rows of $K$ bits. The prover applies the same SPIN encoder to
each row, obtaining $N=2K$ columns, and commits to these columns with a
Merkle tree. An opening sends one random testing combination of the
message rows and one combination weighted by the requested evaluation
point. The combinations lie in $\mathbb F^K$. The verifier encodes them
and checks consistency with authenticated columns sampled after both
messages are fixed. The evaluation combination also supplies the claimed
polynomial value. We retain Brakedown's matrix-testing structure and send
the folded messages explicitly.

\paragraph{Making the binary code useful.}
Expanding every bit to an extension-field element would erase much of
SPIN's encoding advantage. We therefore keep the matrix packed and
encode bounded groups of rows. The internal layout follows the input
index mapping, avoiding a full-matrix bit transpose. Batched column
processing and table-based row combinations reduce the data movement
and arithmetic around the encoder. Workspaces are allocated before
repeated proving. These changes make encoding and column hashing account
for approximately $79\%$ of the fastest standalone prover time.

\paragraph{Parameters.}
We condition once on the selected code satisfying the distance bound
of Theorem~\ref{thm:finite-bch-spin}. Extending its binary generator to
$\mathbb F$ preserves minimum distance. At $K=2^{16}$ and $2^{18}$,
we use $2045$ column samples and one random testing fold. The generic-code
proximity calculation~\cite{diamondPosen24} then gives an error below
$2^{-100}$ for the fixed-matrix test and evaluation check. We call this
the conditional 100-bit parameter target. It excludes the one-time setup
event and the separate commitment-extraction and Fiat--Shamir reductions;
the calculation does not establish a complete composed security theorem.
The artifact records the bound and checks the parameters with exact arithmetic.

\paragraph{Standalone performance.}
Table~\ref{tab:spin-pcs-standalone} compares SPIN--Brakedown with
Ligerito~\cite{ligerito25} on the same $2^{32}$-bit input. SPIN uses
$R=K=2^{16}$ for the square shape. The longer-row shape has $R=2^{14}$
and $K=2^{18}$. Both encode to $1$\,GiB.
Longer rows reduce the SPIN proof
size by $31.7\%$ for a $5.5\%$ increase in prover time. The square
configuration spends $233$\,ms encoding, $179$\,ms hashing columns,
and approximately $95$\,ms computing the two row combinations.

\begin{table}[tb]
\centering
\caption{Standalone PCS for $512$\,MiB of Boolean input. Prover time
includes commitment and one evaluation opening. SPIN uses the conditional
100-bit target; Ligerito uses its Fast100 and Slim100 profiles.}
\label{tab:spin-pcs-standalone}
\small
\begin{tabular}{@{}lrrr@{}}
\toprule
Configuration & Prover (ms) & Verify (ms) & Proof (MiB) \\
\midrule
% Generated by build_application_tables.py; data SHA-256 27480f2ae253a94829a8b2184e483405bdaf91be3971e18ffaf924e23df35ce3
% Ligerito data SHA-256 ddf993619da6426cae615e196bc8485742539c188b7def0bf368feb94a181839
SPIN--Brakedown (square) & 520 & 15 & 18.2 \\
SPIN--Brakedown (longer rows) & 548 & 13 & 12.4 \\
Ligerito (Fast100) & 4331 & 2 & 0.5 \\
Ligerito (Slim100) & 5567 & 1 & 0.3 \\
\bottomrule

\end{tabular}
\end{table}

All measurements use one worker on the Ryzen 9 7950X, pinned to CPU~15,
with a requested $4.5$\,GHz frequency and boost disabled. Each SPIN entry is
the pooled median of ten trials from two processes, each excluding one
warmup. For SPIN, setup, workspace allocation, outer serialization, and independent
evaluation for validation are outside the prover timer. Verification
excludes decoding and workspace allocation. Every timed proof verifies.

Ligerito's rows use the optimized implementation that also supplies the
Flock baseline in Section~\ref{sec:spin-flock}. Its ring-switching layer
opens the binary multilinear polynomial at an ordinary evaluation point;
no Flock computation or cached partial evaluation is supplied. Fast100
and Slim100 use initial code rates $1/2$ and $1/4$, respectively, and
their shipped query and grinding parameters. Their timings pool ten trials
from two processes per profile, each with one excluded warmup. Input
preparation, statement serialization, independent evaluation, outer proof
serialization, and decoding are excluded; internal allocations remain
timed. These profiles give much smaller proofs and faster verification,
at higher prover cost. The artifact records the configurations and the
different timing boundaries; these runs do not establish a common
composed security theorem.

\paragraph{Comparison with Bolt.}
At the same $512$\,MiB input volume and on one core, the fastest measured
Bolt-max commitment takes $1720$\,ms, compared with $413$\,ms for
SPIN--Brakedown: a $4.16\times$ difference. Bolt~\cite{bolt26} encodes
$2^{27}$ elements of $\F_{2^{32}}$; our input has $2^{32}$ Boolean entries.
The comparison therefore matches input bytes, rather than polynomial
dimension or coefficient field. It includes each implementation's code,
layout, and commitment processing. The authors' public Bolt implementation
provides commitment but does not implement the complete opening prover.
The artifact records the
implementations, timing boundaries, and both measured Bolt-max variants.

\paragraph{Opening cost in isolation.}
Table~\ref{tab:pcs-opening} separates opening from commitment. The two
SPIN configurations both open in approximately $107$\,ms, including the
row combinations and authenticated column queries. These are measured
opening times from the same trials as Table~\ref{tab:spin-pcs-standalone};
they exclude outer serialization. Sending the folded messages explicitly
keeps the prover's opening work small, at the proof sizes reported above.

For Bolt, we implemented and measured the dominant opening computations
to obtain a preliminary cost estimate. These estimates indicated a
substantial gap, so we did not pursue a complete Bolt opening implementation
for this comparison. We report a calibrated projection instead.

\begin{table}[tb]
\centering
\caption{One opening after commitment, at $512$\,MiB input volume on one
core. Bolt entries are calibrated projections, not complete prover
measurements or runtime bounds.}
\label{tab:pcs-opening}
\small
\begin{tabular}{@{}llr@{}}
\toprule
Configuration & Evidence & Opening (ms) \\
\midrule
% Generated by build_application_tables.py; data SHA-256 27480f2ae253a94829a8b2184e483405bdaf91be3971e18ffaf924e23df35ce3
% Bolt projection SHA-256 5f80161be50bbea650c1ab082c2b6b79ef05deba1c4b25bb960dfcdff171d22f
SPIN--Brakedown (square) & Measured & 106.82 \\
SPIN--Brakedown (longer rows) & Measured & 107.04 \\
Bolt-max (amortized limit) & Calibrated projection & 1273 \\
Bolt-max (non-amortized) & Calibrated projection & 1729 \\
\bottomrule

\end{tabular}
\end{table}

Bolt's opening additionally uses inner Reed--Solomon proofs and a proof
of the sparse-matrix relation~\cite{bolt26}. To estimate its cost, we
measure row evaluation, a random column fold, and sumchecks, and use
complete Ligerito proximity runs as proxies for the inner proofs. The
calibrated projection is $1273$\,ms in the amortized limit, which assigns
negligible cost per opening to the matrix-relation proof. Adding a
calibration of that proof's leading multiplication work gives $1729$\,ms
without amortization. The inner proxies differ from Bolt's constrained
and batched proofs; basis preparation, some matrix-proof work, outer
authentication, and transcript integration remain unmodeled. The
artifact records these choices and the measured components. Consequently,
these figures describe the projected cost of the selected implementation
approach, rather than a lower bound on Bolt's opening time.

\paragraph{Compatibility with Blaze.}
SPIN is also compatible with Blaze's generic code-switching
framework~\cite{blaze25}. Given the stated distance bound, its linear
encoder can be combined with a generic interactive oracle proof of
proximity for the encoding and multilinear-evaluation relation. This
gives a theoretical SPIN--Blaze composition with asymptotically smaller
proofs than our Brakedown-style implementation. The inner proof replaces
the explicit folded messages, allowing longer rows and shorter
authenticated columns. The large proofs in our measurements are therefore
a property of the chosen PCS construction, not a limitation of SPIN.
We make this compatibility observation without giving a concrete
SPIN--Blaze instantiation, parameter selection, or performance claim.

The opening comparison motivates testing the complete application:
commitment speed alone does not determine the cost of a PCS inside a
proof system. Section~\ref{sec:spin-flock} measures both commitment and
opening after integration into Flock, where the claim interface and
available intermediate computations also affect their cost.
